\documentclass[a4paper,10pt,BCOR10mm,oneside,headsepline]{scrartcl}
\usepackage[portuges]{babel}
\usepackage[utf8]{inputenc}
\usepackage{amsmath}
\usepackage{wasysym}% provides \ocircle and \Box
\usepackage{enumitem}% easy control of topsep and leftmargin for lists
\usepackage{color}% used for background color
\usepackage{forloop}% used for \Qrating and \Qlines
\usepackage{ifthen}% used for \Qitem and \QItem
\usepackage{typearea}
\usepackage{graphicx}
\usepackage{booktabs}
\usepackage{listings}
\usepackage{color}
\usepackage{algorithm}
\usepackage{algorithmic}
\usepackage{amssymb}
\usepackage{minted}
\usepackage{mdframed}
\usepackage{wasysym}
\usepackage{xcolor}

\usepackage{tgtermes}

\areaset{17cm}{26cm}
\setlength{\topmargin}{-2cm}

\definecolor{codegreen}{rgb}{0,0.6,0}
\definecolor{codegray}{rgb}{0.5,0.5,0.5}
\definecolor{codepurple}{rgb}{0.58,0,0.82}
\definecolor{backcolour}{rgb}{0.95,0.95,0.92}
 
\lstdefinestyle{mystyle}{
    backgroundcolor=\color{backcolour},   
    commentstyle=\color{codegreen},
    keywordstyle=\color{magenta},
    numberstyle=\tiny\color{codegray},
    stringstyle=\color{codepurple},
    basicstyle=\footnotesize,
    breakatwhitespace=false,         
    breaklines=true,                 
    captionpos=b,                    
    keepspaces=true,                 
    numbers=left,                    
    numbersep=5pt,                  
    showspaces=false,                
    showstringspaces=false,
    showtabs=false,                  
    tabsize=2
}
 
\lstset{style=mystyle}

\mdfsetup{
    skipabove=12pt,
    skipbelow=12pt,
    innertopmargin=8pt,
    innerbottommargin=8pt,
    innerleftmargin=10pt,
    innerrightmargin=10pt,
    splittopskip=12pt,
    splitbottomskip=12pt,
}

\algsetup{linenodelimiter=.,linenosize=\tiny}

\renewcommand{\algorithmicrequire}{\textbf{Entrada:}}
\renewcommand{\algorithmicensure}{\textbf{Saída:}}
\floatname{algorithm}{Algoritmo}

%% \Qq = Questionaire question. Oh, this is just too simple. It helps
%% making it easy to globally change the appearance of questions.
\newcommand{\Qq}[1]{\textbf{#1}}
%\newcommand{\Qq}[1]{{#1}}

%% \QO = Circle or box to be ticked. Used both by direct call and by
%% \Qrating and \Qlist.
\newcommand{\QO}{$\Box$}% or: $\ocircle$

%% \Qrating = Automatically create a rating scale with NUM steps, like
%% this: 0--0--0--0--0.
\newcounter{qr}
\newcommand{\Qrating}[1]{\QO\forloop{qr}{1}{\value{qr} < #1}{---\QO}}

%% \Qline = Again, this is very simple. It helps setting the line
%% thickness globally. Used both by direct call and by \Qlines.
\newcommand{\Qline}[1]{\noindent\rule{#1}{0.6pt}}

%% \Qlines = Insert NUM lines with width=\linewith. You can change the
%% \vskip value to adjust the spacing.
\newcounter{ql}
\newcommand{\Qlines}[1]{\forloop{ql}{0}{\value{ql}<#1}{\vskip0em\Qline{\linewidth}}}

%% \Qlist = This is an environment very similar to itemize but with
%% \QO in front of each list item. Useful for classical multiple
%% choice. Change leftmargin and topsep accourding to your taste.
\newenvironment{Qlist}{%
\renewcommand{\labelitemi}{\QO}
\begin{itemize}[leftmargin=1.5em,topsep=-.5em]
}{%
\end{itemize}
}

%% \Qtab = A "tabulator simulation". The first argument is the
%% distance from the left margin. The second argument is content which
%% is indented within the current row.
\newlength{\qt}
\newcommand{\Qtab}[2]{
\setlength{\qt}{\linewidth}
\addtolength{\qt}{-#1}
\hfill\parbox[t]{\qt}{\raggedright #2}
}

%% \Qitem = Item with automatic numbering. The first optional argument
%% can be used to create sub-items like 2a, 2b, 2c, ... The item
%% number is increased if the first argument is omitted or equals 'a'.
%% You will have to adjust this if you prefer a different numbering
%% scheme. Adjust topsep and leftmargin as needed.
\newcounter{itemnummer}
\newcommand{\Qitem}[2][]{% #1 optional, #2 notwendig
\ifthenelse{\equal{#1}{}}{\stepcounter{itemnummer}}{}
\ifthenelse{\equal{#1}{a}}{\stepcounter{itemnummer}}{}
\begin{enumerate}[topsep=2pt,leftmargin=2.8em]
\item[\textbf{\arabic{itemnummer}#1.}] #2
\end{enumerate}
}

%% \QItem = Like \Qitem but with alternating background color. This
%% might be error prone as I hard-coded some lengths (-5.25pt and
%% -3pt)! I do not yet understand why I need them.
\definecolor{bgodd}{rgb}{0.8,0.8,0.8}
\definecolor{bgeven}{rgb}{0.9,0.9,0.9}
\newcounter{itemoddeven}
\newlength{\gb}
\newcommand{\QItem}[2][]{% #1 optional, #2 notwendig
\setlength{\gb}{\linewidth}
\addtolength{\gb}{-5.25pt}
\ifthenelse{\equal{\value{itemoddeven}}{0}}{%
\noindent\colorbox{bgeven}{\hskip-3pt\begin{minipage}{\gb}\Qitem[#1]{#2}\end{minipage}}%
\stepcounter{itemoddeven}%
}{%
\noindent\colorbox{bgodd}{\hskip-3pt\begin{minipage}{\gb}\Qitem[#1]{#2}\end{minipage}}%
\setcounter{itemoddeven}{0}%
}
}

\newcommand{\rawcode}[1]{\texttt{#1}}
\newcommand{\code}[2]{\mintinline{#1}{#2}}
\newcommand{\file}[1]{\texttt{\underline{#1}}}

\renewcommand{\theFancyVerbLine}{\tiny \texttt{\textcolor{gray}{\arabic{FancyVerbLine}}}}

\newcommand{\inputcode}[2]{\begin{mdframed}\inputminted[fontsize=\footnotesize,linenos,numbersep=4pt]{#1}{#2}\end{mdframed}}

\newenvironment{solution}{%
    \begin{mdframed}%
        \noindent\textbf{Solu\c{c}\~ao:}\par\medskip
}{%
    \end{mdframed}%
}

\newcommand{\inputsnippet}[4]{\begin{mdframed}[shadow=true,shadowsize=3pt,shadowcolor=gray]\inputminted[fontsize=\footnotesize,linenos,firstline=#2,lastline=#3,numbersep=2pt]{#1}{#4}\end{mdframed}}

\newcommand{\inputsyntax}[2]{\inputminted[fontsize=\footnotesize]{#1}{#2}}



\begin{document}

\begin{center}
	\includegraphics[scale=0.45]{fga}
	% fga.ps: 1087x103 pixel, 72dpi, 38.35x3.63 cm, bb=0 0 1087 103
\end{center}


\vspace{0.1in}

\begin{center}
    \textbf{\huge Paradigmas de Programação}

    \vspace{0.1in}
    \textbf{\Large Programação Funcional: Lista 2}
\end{center}\vskip1em

\renewcommand{\Qq}[1]{{#1}}

\section*{Parte A - Funções de alta ordem}

\Qitem{ \Qq{Usando apenas as funções das bibliotecas \code{haskell}{Prelude} e 
        \code{haskell}{Data.Char} e funções de alta ordem, implemente, de forma tácita,
        as funções listadas abaixo.

    \begin{Qlist}
        \item[(a)] \code{haskell}{count :: String -> Int}, que recebe uma string e retorne o número
            de palavras que terminem com letra minúscula.

                \inputsyntax{haskell}{codes/count.hs}

        \begin{solution}
        \code{haskell}{words} separa a string em palavras, \code{haskell}{last} extrai o último
        caractere de cada palavra (nunca é vazia) e \code{haskell}{isLower} testa se ele é
        minúsculo. \code{haskell}{filter} mantém as palavras que passam no teste e
        \code{haskell}{length} as conta.

        \inputsnippet{haskell}{3}{5}{codes/sol/solutions.hs}


        \textbf{Soluções alternativas:} a primeira usa \code{haskell}{map} para obter uma lista de \code{haskell}{Bool} (um por palavra) e conta os verdadeiros com \code{haskell}{filter id}; a segunda extrai o último caractere de cada palavra com \code{haskell}{map last} e filtra os caracteres minúsculos com \code{haskell}{isLower}.

        \inputsnippet{haskell}{4}{8}{codes/sol/alternatives.hs}
        \end{solution}

        \item[(b)] \code{haskell}{inverses :: [Int] -> [Double]}, que recebe uma lista de inteiros e 
            retorne uma lista com os inversos multiplicativos dos elementos não-nulos. Ex

                \inputsyntax{haskell}{codes/inverses.hs}

        \begin{solution}
        \code{haskell}{filter (/= 0)} descarta os zeros; em seguida, \code{haskell}{fromIntegral}
        converte cada \code{haskell}{Int} em \code{haskell}{Double} e \code{haskell}{(1 /)} calcula
        o inverso.

        \inputsnippet{haskell}{7}{9}{codes/sol/solutions.hs}

        \textbf{Soluções alternativas:} \code{haskell}{recip x} é o mesmo que \code{haskell}{1 / x}.

        \inputsnippet{haskell}{11}{12}{codes/sol/alternatives.hs}
        \end{solution}

        \item[(c)] \code{haskell}{odds :: Int -> Int}, que recebe um inteiro e 
            returne o número de dígitos ímpares deste número.

                \inputsyntax{haskell}{codes/odds.hs}

        \begin{solution}
        \code{haskell}{abs} descarta o sinal, \code{haskell}{show} converte o número em uma string
        de dígitos e \code{haskell}{digitToInt} converte cada caractere de volta em inteiro, que é
        testado com \code{haskell}{odd}.

        \inputsnippet{haskell}{11}{13}{codes/sol/solutions.hs}

        \textbf{Soluções alternativas:} a primeira versão converte cada caractere em inteiro com \code{haskell}{read [c]}; a segunda só testa se o caractere pertence a \code{haskell}{"13579"}.

        \inputsnippet{haskell}{15}{19}{codes/sol/alternatives.hs}
        \end{solution}

        \item[(d)] \code{haskell}{palindromes :: String -> Int}, que recebe uma string e retorna
            o número de palavras que são palíndromos. Ex.:

                \inputsyntax{haskell}{codes/palindromes.hs}

        \begin{solution}
        Para funções, \code{haskell}{f <*> g} aplicada a um argumento $w$ resulta em
        \code{haskell}{f w (g w)}. Logo, \code{haskell}{(==) <*> reverse} aplicada a $w$ resulta em
        \code{haskell}{w == reverse w}, isto é, testa se uma palavra é igual à sua inversa.

        A comparação diferencia maiúsculas de minúsculas, como no exemplo do enunciado:
        \code{haskell}{"Aia"} não é contada, e o resultado é $2$ (``mussum'' e ``mirim'').

        \inputsnippet{haskell}{15}{17}{codes/sol/solutions.hs}

        \textbf{Soluções alternativas:} a gente sempre tem que associar cada palavra o par formado por ela e pela sua inversa, e testar se as duas componentes são iguais. 
        
        Com algo como \code{haskell}{uncurry (==)}, a primeira usa função anônima; a segunda \code{haskell}{(,) <*> reverse}, que aplicada a $w$ resulta em \code{haskell}{(w, reverse w)}; a terceira \code{haskell}{liftA2 (,) id reverse}.

        \inputsnippet{haskell}{22}{29}{codes/sol/alternatives.hs}
        \end{solution}

        \item[(e)] \code{haskell}{isOctNumber :: Int -> Bool}, que recebe um inteiro
            \code{haskell}{x} e retorna verdadeiro se $x$ representa um número em base octal. Ex.:

                \inputsyntax{haskell}{codes/oct.hs}

        \begin{solution}
        Um número está escrito em base octal se todos os seus dígitos estão em
        \code{haskell}{"01234567"}. Assim, \code{haskell}{show . abs} produz a string dos dígitos e
        \code{haskell}{all (flip elem "01234567")} verifica que todos pertencem ao conjunto.

        \inputsnippet{haskell}{19}{21}{codes/sol/solutions.hs}

        \textbf{Soluções alternativas:} a primeira versão define a função auxiliar \code{haskell}{within} (pertinência em lista, com os argumentos invertidos) e a usa no lugar de \code{haskell}{flip elem}; a segunda usa o fato de que os caracteres \code{haskell}{'0'} a \code{haskell}{'9'} são consecutivos e ordenados, de modo que um dígito é octal se, e somente se, é menor que \code{haskell}{'8'}.

        \inputsnippet{haskell}{32}{39}{codes/sol/alternatives.hs}
        \end{solution}
    \end{Qlist}
} }

\Qitem{ \Qq{Usando uma única dobra à esquerda, implemente as funções abaixo. A implementação deve ser
        feita em uma única linha.

    \begin{Qlist}
        \item[(a)] \code{haskell}{eval :: [Int] -> Int -> Int}, que recebe uma lista de inteiros 
            \code{haskell}{as} e um inteiro \code{haskell}{x} e retorna o valor de $p(x)$, onde
            \[
                p(x) = a_0 + a_1x + a_2x^2 + \ldots + a_Nx^N
            \]

        \textbf{Observação}: a lista \code{haskell}{as} contém os coeficientes do polinômio do maior
        índice para o menor índice. Ex.:
                \inputsyntax{haskell}{codes/eval.hs}

        \textbf{Dica:} utilize o algoritmo de Horner. 

        \begin{solution}
        O algoritmo de Horner reescreve
        $p(x) = (\ldots((a_Nx + a_{N-1})x + a_{N-2})x + \ldots)x + a_0$. Como a lista já vem do maior
        para o menor índice, basta percorrê-la da esquerda para a direita com acumulador inicial $0$,
        fazendo $\text{acc} \leftarrow \text{acc}\cdot x + a$ a cada coeficiente $a$.

        \inputsnippet{haskell}{23}{25}{codes/sol/solutions.hs}

        Exemplo ($p(x)=x^2-5x+6$, $x=2$): $((0\cdot 2+1)\cdot 2-5)\cdot 2+6 = 0$.


        \textbf{Soluções alternativas:} \code{haskell}{(+) . (* x)} é a função \code{haskell}{\acc a -> acc * x + a}, escrita de forma tácita: primeiro multiplica o acumulador por $x$ e depois soma o coeficiente.

        \inputsnippet{haskell}{42}{43}{codes/sol/alternatives.hs}
        \end{solution}

        \item[(b)] \code{haskell}{k_factorial :: Int -> Int -> Int}, que recebe os inteiros
            \code{haskell}{n} e \code{haskell}{k} e retorna o $k$-fatorial $n!_{(k)}$ de $n$, onde
        \[
            n!_{(k)} = \left\{ \begin{array}{ll} n\cdot (n - k)!_{(k)}, & \mbox{se}\ n > 0;\\
                1,& \mbox{se}\ -k < n \leq 0;\\
                0,& \mbox{caso contrário}\end{array}\right.
        \]

        Ex.:
                \inputsyntax{haskell}{codes/kfact.hs}

        \begin{solution}
        Seja $k>0$. Se $n\leq -k$, o resultado é $0$, o que é obtido com a lista
        \code{haskell}{[0]}. Caso contrário, o resultado é o produto de
        $n,\ n-k,\ n-2k,\ \ldots$ enquanto os termos forem positivos: ao chegar a um termo $m\leq 0$,
        necessariamente $-k<m\leq 0$, e o caso base vale $1$. Essa sequência é gerada por
        \code{haskell}{iterate (subtract k) n} e cortada por \code{haskell}{takeWhile (> 0)}; a dobra
        \code{haskell}{foldl (*) 1} calcula o produto.

        \inputsnippet{haskell}{27}{29}{codes/sol/solutions.hs}

        Exemplo: \code{haskell}{k_factorial 6 2} $= 6\cdot 4\cdot 2 = 48$.


        \textbf{Soluções alternativas:} em vez de \code{haskell}{if}, acrescenta-se um elemento $0$ à lista com a compreensão de listas \code{haskell}{[0 | n <= -k]}, que é a lista \code{haskell}{[0]} se $n\leq -k$ e a lista vazia caso contrário. Quando $n\leq -k$, o \code{haskell}{takeWhile} resulta em lista vazia e o produto final é $0$.

        \inputsnippet{haskell}{46}{48}{codes/sol/alternatives.hs}
        \end{solution}

        \item[(c)] \code{haskell}{mean :: [Int] -> Double}, que recebe uma lista de inteiros 
            \code{haskell}{xs} e retorna a média aritmética destes números. Ex.:

                \inputsyntax{haskell}{codes/mean.hs}

        \begin{solution}
        Como só é permitida uma dobra, a soma e a quantidade de elementos são acumuladas
        simultaneamente em um par \code{haskell}{(soma, quantidade)}. Ao final,
        \code{haskell}{uncurry (/)} divide a primeira componente pela segunda. Para a lista vazia,
        o resultado é \code{haskell}{NaN} ($0/0$).

        \inputsnippet{haskell}{31}{33}{codes/sol/solutions.hs}

        \textbf{Soluções alternativas:} \code{haskell}{f <*> g} aplicada a uma lista $xs$ resulta em \code{haskell}{f xs (g xs)}; assim, o termo divide a soma de $xs$ pelo seu comprimento. O resultado é o mesmo, mas esta versão percorre a lista duas vezes: uma pela dobra e outra por \code{haskell}{length}, por isso ela não segue estritamente a restrição de usar uma única dobra.

        \inputsnippet{haskell}{51}{52}{codes/sol/alternatives.hs}
        \end{solution}

        \item[(d)] \code{haskell}{arithmetic_progression_sum :: Int -> Int -> Int -> Int} que recebe o
            número de termos \code{haskell}{n}, o primeiro termo \code{haskell}{a} e a razão
            \code{haskell}{q} de uma progressão aritmética e retorne a soma dos primeiros $n$
            termos desta progressão. Ex.:

                \inputsyntax{haskell}{codes/pa.hs}

        \begin{solution}
        \code{haskell}{iterate (+ q) a} gera a progressão infinita $a, a+q, a+2q, \ldots$;
        \code{haskell}{take n} guarda os $n$ primeiros termos e a dobra \code{haskell}{foldl (+) 0}
        os soma.

        \inputsnippet{haskell}{35}{37}{codes/sol/solutions.hs}

        Exemplo: $1+3+5+\cdots+19=100$.


        \textbf{Soluções alternativas:} a primeira versão escreve \code{haskell}{iterate} com \code{haskell}{flip}; a segunda dobra sobre os índices $0, 1, \ldots, n-1$, somando o termo $a + i\,q$ ao acumulador.

        \inputsnippet{haskell}{55}{59}{codes/sol/alternatives.hs}
        \end{solution}

        \item[(e)] \code{haskell}{exp_approx :: Double -> Double}, que recebe um número \code{haskell}{x}
            e retorna a aproximação da exponencial por série de potências até o 11° termo, isto
            é,
        \[
            e^x \approx 1 + \frac{x}{1!} + \frac{x^2}{2!} + \ldots + \frac{x^{10}}{10!}
        \]
        Ex.:
                \inputsyntax{haskell}{codes/exp.hs}

        \begin{solution}
        O termo de índice $n$ é $x^n/n!$, para $n=0,1,\ldots,10$ (onze termos). Esses termos são
        calculados com \code{haskell}{map} e somados com uma única dobra à esquerda.

        \inputsnippet{haskell}{39}{41}{codes/sol/solutions.hs}

        \textbf{Soluções alternativas:} a primeira versão emparelha $x$ com os expoentes por meio de \code{haskell}{zip}; como o \code{haskell}{zip} é truncado no menor operando, só são usados onze termos. A segunda usa o método de Horner: $e^x\approx 1 + x\bigl(1 + \tfrac{x}{2}\bigl(1 + \tfrac{x}{3}(\ldots(1 + \tfrac{x}{10}))\bigr)\bigr)$, calculado por uma dobra à esquerda com $n=10,9,\ldots,1$ e $\text{acc}\leftarrow 1 + \text{acc}\cdot x/n$.

        \inputsnippet{haskell}{62}{68}{codes/sol/alternatives.hs}
        \end{solution}
    \end{Qlist}
} }

\section*{Parte B - I/O}

\Qitem{ \Qq{Resolva, no Beecrowd, o exercício indicado em Haskell. Para formatar a saída, use a função
        \code{haskell}{printf} do módulo \code{haskell}{Text.Printf}.

    \begin{Qlist}
        \item[(a)] \texttt{BEE 1007} -- Diferença

        \begin{solution}

        \inputsnippet{haskell}{1}{6}{codes/sol/bee1007.hs}
        \end{solution}

        \item[(b)] \texttt{BEE 1008} -- Salário

        \begin{solution}

        \inputsnippet{haskell}{1}{7}{codes/sol/bee1008.hs}
        \end{solution}

        \item[(c)] \texttt{BEE 1011} -- Esfera

        \begin{solution}

        \inputsnippet{haskell}{1}{6}{codes/sol/bee1011.hs}
        \end{solution}

        \item[(d)] \texttt{BEE 1013} -- O Maior

        \begin{solution}

        \inputsnippet{haskell}{1}{6}{codes/sol/bee1013.hs}
        \end{solution}

        \item[(e)] \texttt{BEE 1014} -- Consumo

        \begin{solution}

        \inputsnippet{haskell}{1}{6}{codes/sol/bee1014.hs}
        \end{solution}
    \end{Qlist}
} }


\end{document}
